\section{Looping a Causal Encoder--Decoder}
\label{sec:method}

Unlike a homogeneous dense Transformer—where every layer generates and consumes full-context KV, making loop depth a purely representational choice—layer depth in a CED is directly coupled to both \emph{global KV routing} and \emph{prefill execution} (\S\ref{sec:background}). Combining recurrence with a CED therefore requires jointly resolving two questions: \textbf{where} in the asymmetric stack to place the recurrent core (\S\ref{sec:method-where}), and \textbf{how} to train and cache recurrent states so that multi-round iteration does not inflate resident KV memory (\S\ref{sec:method-train}).

\subsection{Where to Loop: Architectural Placement and Serving Asymmetry}
\label{sec:method-where}

Because a weight-tied recurrent block cannot straddle the encoder--decoder boundary without breaking the single-projection invariant that enables decoder-free prefill, the CED topology naturally partitions loop placement into two regimes (Figure~\ref{fig:main}a):

\textbf{Encoder loop (Iterative Global KV Refinement).} The recurrent core resides entirely within the encoder and encompasses at least one global KV source layer. Across the $T$ rounds, the core iteratively refines both its local sliding-window representations and the encoder's global KV source. Because the final encoder representation—and by extension the projected decoder global KV—depends on the last round, prompt prefill must run all $T$ encoder rounds.

\textbf{Decoder loop (Static-Context Iterative Readout).} The recurrent core resides entirely within the decoder, downstream of the boundary projection. Here, all global KV entries—both in the encoder and at the boundary—are finalized before the loop begins. Each recurrent round refines its local SWA states while repeatedly querying a \emph{static} global KV cache. Because the decoder is skipped during prefill, this placement leaves prefill compute unchanged relative to an unlooped CED, apart from the last $w$ prompt tokens, which run the decoder to fill its SWA caches.

\paragraph{Compute and execution asymmetry.} 
In a CED, representational role and system execution move in lockstep. During autoregressive decoding, every token traverses the full network, so looping a block of layers for $T$ rounds adds the exact same number of single-layer passes per token in either placement. During prompt prefill, however, a decoder loop executes only the unlooped encoder, whereas an encoder loop must execute all $T$ rounds of its recurrent core before projecting the boundary KV. 

In our 12-layer experimental configuration with a 6-layer encoder, a 4-layer recurrent core, and $T{=}\val{cfg/rounds}$ rounds (Fig.~\ref{fig:main}a), an unlooped CED requires \val{passes/x40/1024/ced-noloop/prefill} layer passes per prompt token (plus the boundary projection), halving the \val{passes/x40/1024/dense-noloop/prefill} passes of an unlooped dense baseline. Looping the decoder preserves this lightweight \val{passes/x40/1024/ced-dec/prefill}-pass prefill, whereas looping the encoder expands prefill to \val{passes/x40/1024/ced-enc/prefill} layer passes per token; during decode, both placements expand execution from \val{passes/x40/1024/ced-noloop/decode} to \val{passes/x40/1024/ced-enc/decode} passes. The two placements also differ in KV footprint. Decoder layers produce no global KV, so a decoder loop keeps the unlooped CED's \val{kvlast/x40/1024/ced-noloop}\,KB of global KV per token (150M models, bf16) even if every round caches its own KV; an encoder loop rewrites its global KV source in every round, so per-round caching grows the global KV to \val{kvround/x40/1024/ced-enc}\,KB per token. Thus, comparing encoder and decoder loops directly tests whether high-quality recurrence requires actively refining the global KV source at the cost of prefill compute, or whether iterating over a fixed global context during decode suffices.

% \subsection{How to Loop: Fixed-Point Training and Dual-Pathway KV Sharing}
\subsection{How to Loop: Fixed-Point Training and Last-Round KV Caching}
\label{sec:method-train}

Even when a loop improves modeling quality, naively caching distinct keys and values across all $T$ rounds multiplies the resident KV footprint of every looped layer by $T$, negating CED's memory advantage. Eliminating this factor requires bridging a fundamental discrepancy between parallel training and memory-efficient inference.

\paragraph{Per-round training vs.\ last-round inference.}
During training, sequences are processed in parallel across tokens by unrolling the recurrent core: at round $r$, each layer inside the core attends to the round-$r$ keys and values of preceding tokens, which we term \emph{per-round KV}. Caching all $T$ trajectories during serving would require $T$ copies of the core's global and SWA caches: for the 150M encoder loop at $T{=}\val{cfg/rounds}$, the global KV per token grows from \val{kvlast/x40/1024/ced-enc}\,KB to \val{kvround/x40/1024/ced-enc}\,KB (\S\ref{sec:method-where}; Fig.~\ref{fig:main}b, hatched bars). To avoid this overhead, at inference we apply \emph{last-round KV} (or terminal KV sharing~\citep{loopsright2,benice}): once a token completes all $T$ rounds, we retain only its last-round KV states—both the global KV from the core's source layer and the local SWA KV of every looped layer—and discard all intermediate rounds. When a subsequent token executes any round $r$, it attends to its own round-$r$ state for the current step alongside the cached last-round KV for all past tokens.

\paragraph{Fixed-point convergence under stochastic depth.}
Substituting last-round KV for intermediate-round KV at inference avoids distribution shift only if the recurrent core converges toward a contractive fixed point across rounds, so that both the local SWA representations and the global KV source stabilize in late rounds. Unlike a deterministic fixed-depth unroll—which allows the core to learn round-specific feature transformations that diverge across rounds—we sample a stochastic recurrence depth per training sequence from a Poisson distribution with mean $\val{cfg/rounds}$ (clipped to at least $1$, backpropagating gradients through the final $\val{cfg/bwd}$ rounds) while executing a deterministic $\val{cfg/rounds}$ rounds at inference (Table~\ref{tab:hparams}). Exposing the coda and downstream layers to recurrent states across variable depths pressures the core to settle rapidly into a depth-invariant equilibrium~\citep{loopsright2,benice}. When this equilibrium holds (\S\ref{sec:res-kv}), \ourmethod retains the exact same resident KV footprint per token as an unlooped CED (\val{kvlast/x40/1024/ced-enc}\,KB of global KV per token at 150M, against \val{kvround/x40/1024/ced-enc}\,KB with per-round caching).
